\subsection{Measurement, contextuality, and effective descriptions}
  \label{subsec:contextuality}

  The structural origin of Bell-type correlations identified above naturally raises the
  question of their suppression in classical macroscopic regimes.
  Everyday physical systems exhibit effectively local and statistically independent
  behavior, despite being composed of quantum constituents.

  Recent work in quantum foundations has emphasized that nonlocal and contextual
  correlations may be understood as obstructions to constructing globally consistent
  probabilistic models from local measurement contexts~\cite{AbramskyBrandenburger2011}.
  Within ontological-model frameworks, observable statistics are likewise interpreted
  as coarse-grained descriptions of an underlying ontic state space~\cite{Spekkens2007}.
  The mechanism discussed here is complementary in spirit, but identifies a more
  primitive structural origin for the failure of Bell factorizability: the
  non-injective mapping between underlying configurations and observable
  descriptions.

  Within the present framework, this transition is understood as a change in the
  effective properties of the descriptive mapping.
  Environmental interactions, decoherence, and coarse-graining progressively reduce
  the number of distinct underlying configurations compatible with a given observable
  outcome.
  As a result, the effective mapping becomes approximately injective, and joint
  probabilities recover an approximately factorizable form.

  In this regime, observable outcomes may be treated as effectively independent, and
  Bell inequalities are no longer violated in practice.
  Classical statistical independence thus emerges as a limiting case of effective
  injectivity, rather than as a fundamental principle.

  The present work does not propose a modification of quantum mechanics or a new
  ontological framework.
  Its contribution is to isolate a minimal and sufficient structural condition for the
  failure of Bell inequalities, thereby clarifying the conceptual origin of quantum
  correlations while preserving both operational locality and empirical adequacy.

  We conclude that Bell inequality violations need not be interpreted as evidence for
  fundamental nonlocality, but may instead reflect intrinsic limitations of effective
  descriptions in which observable outcomes are treated as independently specifiable
  local properties.
